\documentclass[10pt, twocolumn]{article}
\usepackage[margin=1.5cm]{geometry}
\usepackage[T1]{fontenc}
\usepackage{lmodern}
\usepackage[protrusion=true,expansion=true,final]{microtype}
\usepackage[hidelinks]{hyperref}
\setlength{\columnsep}{0.6cm}
\setlength{\parindent}{0pt}
\setlength{\parskip}{2pt}

%% Packages
\usepackage{amsmath,amssymb}
\usepackage{graphicx}
\usepackage{xcolor}
\usepackage{booktabs}
\usepackage{array}
\usepackage{makecell}
\usepackage{enumitem}
\usepackage{caption}
\usepackage{float}
\usepackage{xspace}
\usepackage{url}

%% Vancouver numbered citations
\usepackage[style=vancouver, backend=biber]{biblatex}
\addbibresource{references.bib}

%% Commands
\newcommand{\varforge}{\textsc{VarForge}\xspace}

%%═════════════════════════════════════════════════════════════════════════
\begin{document}

\title{\varforge: Synthetic Cancer Sequencing Data in a Single Rust Binary\\[0.2ex]
       \small\texttt{v0.1}}
\author{Trent Zeng}
\date{}
\maketitle

\begin{abstract}
\noindent
No single tool integrates de-novo read generation, somatic variant spike-in, UMI/duplex simulation, cfDNA fragmentation, and library artefact injection.
\varforge{} fills this gap in a single Rust binary driven by YAML configuration, achieving ${\sim}7{,}000$ read pairs/s at 30x on a 12-thread workstation.
No other published tool provides genome-wide UMI simulation or combines cfDNA fragmentation with variant spike-in.
Variants are assigned by per-read Bernoulli sampling, not deterministic spike-in.
BED region filtering enables clinical-depth panel simulation (200x, 1\,Mbp target, 86\,s).
\end{abstract}

%%─────────────────────────────────────────────────────────────────────────
\section{Introduction}

Benchmarking cancer variant callers, UMI deduplication tools, and liquid biopsy algorithms requires synthetic data with complete ground truth.
Real patient data cannot provide this.
Existing tools address parts of the problem: BAMSurgeon\autocite{ewing2015bamsurgeon} spikes variants, ART\autocite{huang2012art} generates reads, but no single tool covers UMI families, cfDNA fragmentation, tumour models, and library artefacts together.

\varforge{} addresses these gaps.
Contributions: (i) unified pipeline via YAML; (ii) genome-wide UMI and duplex simulation; (iii) cfDNA fragmentation with variant support; (iv) stochastic VAF via per-read Bernoulli sampling; (v) BED region filtering for targeted panel simulation.

%%─────────────────────────────────────────────────────────────────────────
\section{Related Work}

BAMSurgeon\autocite{ewing2015bamsurgeon}, SomatoSim\autocite{kuhnemund2017somatomsim}, and stochasticSim\autocite{osullivan2023stochasticsim} spike variants into existing BAMs.
ART\autocite{huang2012art}, NEAT\autocite{stephens2016neat}, and ReSeq\autocite{schmeing2021reseq} generate reads but lack tumour models, UMI simulation, and cfDNA fragmentation.
No existing tool provides genome-wide UMI or duplex sequencing simulation.
No tool combines cfDNA fragmentation with somatic variant spike-in.

%%─────────────────────────────────────────────────────────────────────────
\section{Method}

\varforge{} partitions the reference into chunks processed by Rayon workers.
Each worker samples fragment lengths, positions reads, injects variants by Bernoulli sampling, applies quality scores, and tags UMI families.
A bounded channel (default depth 64) decouples workers from the writer thread.
Memory is bounded by concurrent batch state, not total dataset size.
If a BED file of targets is provided, chunks are clipped to target boundaries before dispatch.

\textbf{Stochastic VAF.}
For each read at a variant site: $X_i \sim \text{Bernoulli}(\text{VAF}_\text{eff})$.
Alt-read count follows $\text{Binomial}(D, \text{VAF}_\text{eff})$, matching real sequencing variance.

\textbf{cfDNA.}
Fragment lengths follow a three-component Gaussian mixture with ctDNA (${\sim}143$\,bp), mono-nucleosomal (${\sim}167$\,bp), and di-nucleosomal (${\sim}334$\,bp) components.

%%─────────────────────────────────────────────────────────────────────────
\section{Results}

\begin{table}[H]
\centering
\caption{Key benchmark results (12 threads, 1 run).}
\small
\setlength{\tabcolsep}{3pt}
\begin{tabular}{@{}lrrr@{}}
\toprule
Scenario & Wall (s) & RAM (GB) & Pairs/s \\
\midrule
WGS 10\,MB, 30x  & 139.6 & 0.83 &  7,200 \\
Panel+UMI 1\,MB  & 195.1 & 1.72 &  6,800 \\
chr22, 30x WGS   & 595.3 & 4.11 &  8,537 \\
Panel BED 200x   &  85.5 & 0.55 &  7,794 \\
\bottomrule
\end{tabular}
\end{table}

Throughput scales linearly with coverage ($R^2 > 0.999$).
Peak memory is bounded by concurrent batch count, not dataset size.
BED filtering restricts simulation to target regions, enabling 200x panel coverage on chr22 in 86\,s.
All six hg38 chr22 feature configurations (UMI, duplex, cfDNA, FFPE) completed without error.

%%─────────────────────────────────────────────────────────────────────────
\section{Discussion and Future Work}

\varforge{} covers a feature space unavailable in any single existing tool.
VAF accuracy is validated by a chi-squared goodness-of-fit test against the theoretical binomial distribution.
cfDNA fragment distributions and UMI family structure are not yet independently verified and require quantitative comparison against their respective models.
Thread scaling only helps for compute-heavy workloads (UMI simulation); the I/O-bound WGS workload runs fastest at one thread.

A four-scenario downstream benchmark (bwa-mem2, Mutect2, rtg-tools vcfeval) confirms that VarForge data exposes real caller behaviour: the SNV baseline achieves F1\,=\,0.857, while FFPE artefacts overwhelm stock Mutect2 with 113K false positives.
Priority next steps: (1) end-to-end VAF validation via alignment and pileup; (2) cfDNA fragment distribution validation; (3) extending the variant caller benchmark to additional tools (VarDict, Strelka2, fgbio).

%%─────────────────────────────────────────────────────────────────────────
\printbibliography[heading=bibintoc, title={References}]

\end{document}
